
\documentclass[11pt]{article}
\usepackage{acl}
\usepackage{times}
\usepackage{latexsym}
\usepackage{graphicx}
\usepackage{booktabs}

\title{Multiclass Classification of Cultural Items}
\author{Your Name \\ Sapienza University of Rome}
\date{}

\begin{document}
\maketitle

\section{Introduction}

Understanding cultural specificity is an important challenge for Natural Language Processing (NLP), particularly in tasks involving multicultural knowledge and bias mitigation. Cultural items may be strongly associated with a specific culture, partially representative of it, or culturally agnostic. The goal of this homework is to develop an automatic classifier that categorizes items into three classes: Cultural Agnostic (CA), Cultural Representative (CR), and Cultural Exclusive (CE).

Following the assignment requirements, we implemented two different approaches:
\begin{itemize}
\item A Transformer-based language model classifier.
\item A non-LM approach based on Word2Vec embeddings.
\end{itemize}

Both models use the item name and description as input features and are evaluated on the provided development set using standard classification metrics.

\section{Methodology}

\subsection{Transformer-based Classifier}

The first approach follows a standard supervised text classification pipeline using a pretrained Transformer encoder. We fine-tune a pretrained language model to classify items into the three cultural categories.

The model takes as input the item name concatenated with the description. The input is tokenized using the model tokenizer and passed through the encoder.

We use the hidden representation of the [CLS] token from the final layer as the aggregated representation of the input sequence. This vector is then passed through a linear classification layer producing logits for the three cultural classes.

The model is trained using cross-entropy loss and optimized with AdamW. Fine-tuning allows the model to adapt pretrained linguistic knowledge to the specific cultural classification task.

\subsection{Word2Vec-based Classifier}

As a non-LM approach, we implemented a classifier based on Word2Vec embeddings.

Each item description is tokenized and mapped to its corresponding embedding vectors. To obtain a sentence-level representation, we compute the masked mean pooling of token embeddings:

\begin{equation}
v_{sent} = \frac{\sum_i m_i \cdot v_i}{\sum_i m_i}
\end{equation}

where $v_i$ is the embedding of token $i$ and $m_i$ is the attention mask.

The resulting sentence vector is then passed through a feed-forward neural classifier composed of:
\begin{itemize}
\item Layer normalization
\item Dropout regularization
\item A hidden linear layer with ReLU activation
\item A final linear classification layer
\end{itemize}

To mitigate dataset imbalance we apply class weights in the loss function based on the label distribution in the training data.

\section{Experiments}

\subsection{Dataset}

The dataset consists of cultural items extracted from Wikidata. Each entry includes:

\begin{itemize}
\item item name
\item description
\item type
\item hypernym
\item classification label
\end{itemize}

Items are annotated according to the taxonomy:
\begin{itemize}
\item Cultural Agnostic (CA)
\item Cultural Representative (CR)
\item Cultural Exclusive (CE)
\end{itemize}

The training set contains silver annotations while the development set contains gold annotations manually curated by the dataset creators.

\subsection{Training Setup}

\textbf{Transformer model}

\begin{itemize}
\item pretrained encoder: Transformer encoder (e.g. BERT-like model)
\item optimizer: AdamW
\item loss: cross-entropy
\item training epochs: 10
\item input: item name + description
\end{itemize}

\textbf{Word2Vec model}

\begin{itemize}
\item embedding dimension: 300
\item pooling: masked mean pooling
\item classifier: two-layer MLP
\item dropout: 0.1
\item loss: weighted cross-entropy
\end{itemize}

\section{Results}

\begin{center}
\begin{tabular}{lcccc}
\toprule
Model & Accuracy & Precision & Recall & F1 \\
\midrule
Word2Vec + MLP & 0.69 & 0.68 & 0.68 & 0.68 \\
Transformer LM & TBD & TBD & TBD & TBD \\
\bottomrule
\end{tabular}
\end{center}

The Word2Vec model achieves moderate performance, reaching approximately 0.68 macro-F1 on the development set.

The Transformer model is expected to outperform the Word2Vec approach due to its contextual representations and ability to capture deeper semantic relations.

\section{Discussion}

The comparison between the two approaches highlights the trade-off between simplicity and expressive power.

The Word2Vec-based model is computationally efficient and relatively easy to train. However, static embeddings cannot capture contextual meaning, which is often crucial for identifying cultural specificity.

In contrast, Transformer models leverage contextual embeddings and large-scale pretraining, enabling them to better interpret nuanced descriptions.

Another challenge arises from class ambiguity. Some items may lie between cultural representative and cultural exclusive categories, which makes the classification task difficult even for humans.

\section{Conclusion}

In this work we explored two approaches for the classification of cultural items:

\begin{itemize}
\item a Transformer-based language model classifier
\item a Word2Vec-based neural classifier
\end{itemize}

The Word2Vec model provides a lightweight baseline while still achieving reasonable performance. However, contextual Transformer models are better suited for capturing cultural semantics in textual descriptions.

Future work could incorporate additional knowledge sources such as Wikipedia text, Wikidata relations, or external cultural knowledge bases.

\end{document}
